% from ega4-18-fr.tex
\oldpage[IV-4]{194}

\begin{definition}[19.3.1]
\EGAReaderTarget{ega4.c04.def.19-3-1}{IV.19.3.1}
Soit $A$ un anneau local noethérien. On dit que $A$ est un \emph{anneau d'intersection complète} (ou encore, \emph{anneau d'intersection complète absolue}, s'il y a danger de confusion) si le complété $\widehat A$ est isomorphe au quotient d'un anneau local noethérien complet régulier $B$ par un idéal régulier (i.e. (16.9.7) un idéal engendré par une suite régulière d'éléments de $B$).

Tout anneau local régulier est un anneau d'intersection complète.

On dit qu'un préschéma localement noethérien $X$ est \emph{intersection complète} (absolue) \emph{au point} $x\in X$ si l'anneau $\sh{O}_{X,x}$ est un anneau d'intersection complète.
\end{definition}

\begin{proposition}[19.3.2]
\label{IV.19.3.2}
Soient $B$ un anneau local noethérien régulier, $\mathfrak J$ un idéal de $B$. Pour que $A=B/\mathfrak J$ soit un anneau d'intersection complète, il faut et il suffit que l'idéal $\mathfrak J$ soit régulier (i.e. (16.9.7) engendré par une suite régulière d'éléments de $B$).
\end{proposition}

On peut écrire $\widehat A=\widehat B/\mathfrak J\widehat B$ et puisque $\widehat B$ est un $B$-module fidèlement plat et $B$ noethérien, $\mathfrak J\widehat B$ est régulier si et seulement si $\mathfrak J$ l'est (19.1.5, (ii)). Ceci montre (en vertu de ($0$, 17.1.5)) d'une part que la condition de l'énoncé est suffisante, et d'autre part que pour prouver qu'elle est nécessaire, on peut se borner au cas où $A$ et $B$ sont complets. Par hypothèse, il existe alors un anneau local noethérien complet régulier $B'$ tel que $A$ soit isomorphe à un quotient $B'/\mathfrak J'$, où $\mathfrak J'$ est un idéal régulier de $B'$. Considérons alors le \emph{produit fibré} $B\times_A B'=B''$ pour les homomorphismes surjectifs $B\to A$, $B'\to A$ ($0$, 18.1.2); nous prouverons d'abord le lemme suivant:

\begin{lemma}[19.3.2.1]
\label{IV.19.3.2.1}
Soient $A$, $E$, $F$ trois anneaux, $f:E\to A$, $g:F\to A$ deux homomorphismes, et posons $G=E\times_A F$ ($0$, 18.1.2).
\begin{enumerate}
  \item[(i)] Si $f$ est surjectif et si $E$ et $F$ sont des anneaux noethériens, il en est de même de $G$.
  \item[(ii)] Si $A$, $E$ et $F$ sont des anneaux locaux et si $f$ et $g$ sont des homomorphismes locaux, $G$ est un anneau local et les homomorphismes canoniques $G\to E$, $G\to F$ sont locaux.
  \item[(iii)] Si les conditions de (i) et (ii) sont vérifiées, si $g$ est surjectif, et si $E$ et $F$ sont des anneaux locaux noethériens complets, il en est de même de $G$.
\end{enumerate}
\end{lemma}

\begin{enumerate}
  \item[(i)] Compte tenu de ($0$, 18.1.3 et 18.1.5), $G$ est un $A$-anneau augmenté sur $F$, dont l'idéal d'augmentation $\mathfrak J'$ s'identifie canoniquement au noyau $\mathfrak J$ de $f$, la structure de $G$-module sur $\mathfrak J'$ s'identifiant à celle de $E$-module sur $\mathfrak J$ au moyen de l'homomorphisme canonique $G\to E$. Si $\mathfrak a$ est un idéal quelconque de $G$, $\mathfrak a/(\mathfrak a\cap\mathfrak J')$ est isomorphe à $(\mathfrak a+\mathfrak J')/\mathfrak J'$, donc à un idéal de $F$, et par suite est de type fini; d'autre part, $\mathfrak a\cap\mathfrak J'$ est isomorphe à un idéal de $E$ contenu dans $\mathfrak J$, donc est aussi de type fini, ce qui prouve que $\mathfrak a$ est de type fini, donc $G$ est noethérien.
  \item[(ii)] Soient $\mathfrak r$, $\mathfrak s$ les idéaux maximaux de $E$ et $F$ respectivement; l'ensemble $\mathfrak m=\mathfrak r\times_A\mathfrak s$ est évidemment un idéal de $G$, image réciproque de $\mathfrak r$ par la projection $G\to E$ et de $\mathfrak s$ par la projection $G\to F$. On en déduit aussitôt que les éléments de $\mathfrak m$ sont non inversibles dans $G$. D'autre part, si $(x,y)\notin\mathfrak m$, et si par exemple $x\notin\mathfrak r$, $f(x)$ n'appartient pas à l'idéal maximal de $A$, donc, puisque $g(y)=f(x)$, et que $g$ est local, on a aussi $y\notin\mathfrak s$; on en

